\section{Conclusion}
\label{sec:reevaluation-scripts-conclusion}

Dans ce chapitre, nous proposons d'étendre notre mécanisme de réévaluation afin
de supporter les opérations construites à l'aide du langage de scripts de
Jerboa. %
Nous étudions plus particulièrement trois structures de contrôles qui sont
respectivement les séquences, les boucles \texttt{foreach} et les alternatives
\texttt{if-then-else}. %
Une première problématique que nous abordons concerne le mode de représentation
des scripts pour déterminer comment traiter leurs contenus. %
La première option concernait la représentation boîte fermée. %
Dans ce cas, un script ne nécessite pas de traitement particulier et, en
contrepartie, les développeurs de scripts doivent construire une sortie à la
manière des règles Jerboa. %
La seconde option, que nous avons choisie, est la boîte ouverte. %
Dans ce mode de représentation, nous avons connaissances du contenu des scripts,
à savoir les règles appelées et leurs paramètres ainsi que les structures de
contrôle utilisées. %
Dans cette démarche, nous proposons d'étendre nos structures de spécification
paramétrique et de DAGs. %
L'extension des spécifications paramétriques consiste à représenter un script
comme une spécification imbriquée et à lister son contenu, c'est-à-dire les
opérations appliquées et les structures de contrôles utilisées. %
Pour les DAGs, nous considérons qu'il est nécessaire d'ajouter une branche pour
chacun des paramètres qui conditionnent la réévaluation d'une orbite. %
Ainsi, dont l'évolution dépend d'un script, est conditionnée par l'existance du
paramètre topologique de ce script. %
De même, nous considérons chaque itération d'une boucle \texttt{foreach} dispose
de deux paramètres topologiques qui sont l'orbite à itérer et l'instance
d'itération qui impacte l'orbite que nous souhaitons apparier. %
Enfin, nous proposons une dernière contribution relative à la structure
d'alternative \texttt{if-then-else}. %
Cette structure permet, étant donné l'évaluation d'une condition, appliquer une
séquence d'instruction ou une autre. %
Dans le cas où une autre branche que celle évaluée initialement serait
appliquée, il est nécessaire d'identifier précisément les orbites sur lesquelles
les utilisateurs souhaitent travailler. %
En effet deux orbites appelées de la sorte peuvent être très différentes. %
Pour respecter les intentions de modélisation, nous procédons une procédure,
basée sur l'analyse des règles alternatives, qui fait l'appariement des orbites
qui subissent des transformations similaires. %
Enfin, puisque ces appariement ne sont explicitement exprimés par les
utilisateurs, nous proposons la sélection de stratégies pour configurer la
qualité d'appariement recherchée. %
